\documentclass[12pt,letterpaper]{article}
\usepackage{graphicx,textcomp}
\usepackage{natbib}
\usepackage{float}
\usepackage{setspace}
\usepackage{fullpage}
\usepackage{color}
\usepackage[reqno]{amsmath}
\usepackage{amsthm}
\usepackage{graphicx}
\usepackage{fancyvrb}
\usepackage{amssymb,enumerate}
\usepackage[all]{xy}
\usepackage{endnotes}
\usepackage{lscape}
\newtheorem{com}{Comment}
\usepackage{float}
\usepackage{hyperref}
\newtheorem{lem} {Lemma}
\newtheorem{prop}{Proposition}
\newtheorem{thm}{Theorem}
\newtheorem{defn}{Definition}
\newtheorem{cor}{Corollary}
\newtheorem{obs}{Observation}
\usepackage[compact]{titlesec}
\usepackage{dcolumn}
\usepackage{pdfpages}
\usepackage{tikz}
\usetikzlibrary{arrows}
\usepackage{multirow}
\usepackage{xcolor}
\newcolumntype{.}{D{.}{.}{-1}}
\newcolumntype{d}[1]{D{.}{.}{#1}}
\definecolor{light-gray}{gray}{0.65}
\usepackage{url}
\usepackage{listings}
\usepackage{color}

\definecolor{codegreen}{rgb}{0,0.6,0}
\definecolor{codegray}{rgb}{0.5,0.5,0.5}
\definecolor{codepurple}{rgb}{0.58,0,0.82}
\definecolor{backcolour}{rgb}{0.95,0.95,0.92}

\lstdefinestyle{mystyle}{
	backgroundcolor=\color{backcolour},   
	commentstyle=\color{codegreen},
	keywordstyle=\color{magenta},
	numberstyle=\tiny\color{codegray},
	stringstyle=\color{codepurple},
	basicstyle=\footnotesize,
	breakatwhitespace=false,         
	breaklines=true,                 
	captionpos=b,                    
	keepspaces=true,                 
	numbers=left,                    
	numbersep=5pt,                  
	showspaces=false,                
	showstringspaces=false,
	showtabs=false,                  
	tabsize=2
}
\lstset{style=mystyle}
\newcommand{\Sref}[1]{Section~\ref{#1}}
\newtheorem{hyp}{Hypothesis}

\title{Problem Set 1}
\date{Due: October 8, 2026}
\author{
	Applied Stats/Quant Methods 1\\
	Kanon Saito
	}

\begin{document}
	\maketitle
	
	\section*{Instructions}
	\begin{itemize}
	\item Please show your work! You may lose points by simply writing in the answer. If the problem requires you to execute commands in \texttt{R}, please include the code you used to get your answers. Please also include the \texttt{.R} file that contains your code. If you are not sure if work needs to be shown for a particular problem, please ask.
\item Your homework should be submitted electronically on GitHub.
\item This problem set is due before 23:59 on Thursday October 8, 2026. No late assignments will be accepted.
	\end{itemize}
	
	\vspace{1cm}
	\section*{Question 1: Education}

A school counselor was curious about the average of IQ of the students in her school and took a random sample of 25 students' IQ scores. The following is the data set:\\
\vspace{.5cm}

\lstinputlisting[language=R, firstline=36, lastline=36]{PS01_answers_KS.R}  

\vspace{.5cm}

\begin{enumerate}
	\item Find a 90\% confidence interval for the average student IQ in the school (As a hint, you first need the mean and SD to create a CI):\\
	R:
	\lstinputlisting[firstline=37, lastline=42]{PS01_answers_KS.R} 
	Answer: [93.960, 102.92]
	\item Next, the school counselor was curious  whether  the average student IQ in her school is higher than the average IQ score (100) among all the schools in the country.\\ 
	\noindent Using the same sample, conduct the appropriate hypothesis test with $\alpha=0.05$.
	\[
	H_0: \mu \le 100
	\]
	\[
	H_1: \mu > 100
	\]
	R:
	\lstinputlisting[firstline=45, lastline=46]{PS01_answers_KS.R}
	Test Statistic (t-value): -0.596\\
	df: 24\\
	p-value: 0.722\\
	
	Answer: \\
	Since the p-value is greater than 0.05, the null hypothesis cannot be rejected at the 5\% significance level. Therefore, there is not enough statistical evidence to conclude that the average student IQ in this school is higher than the national average.
\end{enumerate}

\newpage

	\section*{Question 2: Political Economy}

\noindent Researchers are curious about what affects the amount of money communities spend on addressing homelessness. The following variables constitute our data set about social welfare expenditures in the USA. \\
\vspace{.5cm}


\begin{tabular}{r|l}
	\texttt{State} &\emph{50 states in US} \\
	\texttt{Y} & \emph{per capita expenditure on shelters/housing assistance in state}\\
	\texttt{X1} &\emph{per capita personal income in state} \\
	\texttt{X2} &  \emph{Number of residents per 100,000 that are "financially insecure" in state}\\
	\texttt{X3} &  \emph{Number of people per thousand residing in urban areas in state} \\
	\texttt{Region} &  \emph{1=Northeast, 2= North Central, 3= South, 4=West} \\
\end{tabular}

\vspace{.5cm}
\noindent Explore the \texttt{expenditure} data set and import data into \texttt{R}.
\vspace{.5cm}
\lstinputlisting[language=R, firstline=42, lastline=42]{PS01_answers_KS.R}  
\vspace{.5cm}
\begin{itemize}

\item
Please plot the relationships among \emph{Y}, \emph{X1}, \emph{X2}, and \emph{X3}? What are the correlations among them (you just need to describe the graph and the relationships among them)?

R:
\lstinputlisting[firstline=55, lastline=67]{PS01_answers_KS.R}
Result: \\
Correlation Matrix
\begin{verbatim}
           Y        X1        X2        X3
Y  1.0000000 0.5317212 0.4482876 0.4636787
X1 0.5317212 1.0000000 0.2056101 0.5952504
X2 0.4482876 0.2056101 1.0000000 0.2210149
X3 0.4636787 0.5952504 0.2210149 1.0000000
\end{verbatim}
\begin{figure}[ht]
\centering
\vspace{.5cm}
\includegraphics[width=10cm]{Rplot_X1_Y.pdf}\\
\vspace{.5cm}
\parbox{15cm}{As the per capita expenditure on shelters/housing assistance (Y) increases, per capita personal income (X1) also tends to increase. The correlation coefficient (0.531) indicates a moderate positive correlation between Y and X1.}
\end{figure}

\begin{figure}[ht]
\centering
\vspace{.5cm}
\includegraphics[width=10cm]{Rplot_X2_Y.pdf}

\vspace{1cm}
\parbox{15cm}{When the number of financially insecure residents per 100,000 (X2) is below about 300, per capita expenditure on shelters/housing assistance (Y) tends to decrease as X2 increases. However, when X2 is above 300, Y tends to increase as X2 increases. The correlation coefficient (0.448) indicates a moderate positive correlation between Y and X2 overall.}
\end{figure}

\begin{figure}[H]
	\centering
	\vspace{.5cm}
	\includegraphics[width=10cm]{Rplot_X3_Y.pdf}

	\parbox{15cm}{As the per capita expenditure on shelters/housing assistance (Y) increases, number of residents per thousand residing in urban areas (X3) also tends to increase. The correlation coefficient (0.448) indicates a moderate positive correlation between Y and X3.}
\end{figure}
\item
Please plot the relationship between \emph{Y} and \emph{Region}? On average, which region has the highest per capita expenditure on housing assistance?\\

R:
\lstinputlisting[firstline=71, lastline=79]{PS01_answers_KS.R}
\vspace{.5cm}
\begin{figure}[H]
	\centering
	\vspace{.5cm}
	\includegraphics[width=10cm]{Rplot_boxplot.pdf}\\
	\vspace{.5cm}
	\parbox{15cm}{The graph shows that the West region has the highest average per capita expenditure on shelters/housing assistance.}
\end{figure}
\item
Please plot the relationship between \emph{Y} and \emph{X1}? Describe this graph and the relationship. Reproduce the above graph including one more variable \emph{Region} and display different regions with different types of symbols and colors.\\

R:
\lstinputlisting[firstline=83, lastline=86]{PS01_answers_KS.R}
\begin{figure}[H]
	\centering
	\vspace{.5cm}
	\includegraphics[width=12cm]{Rplot_X1_Y_region.pdf}\\
	\vspace{.5cm}
	\parbox{15cm}{The graph suggests that there are positive correlations between per capita expenditure (Y) and personal income (X1) in the South and Northeast regions, while the relationship between X1 and Y seems less clear in the North Central and West regions.}
\end{figure}
\end{itemize}


\end{document}
